\subsection{The actual infinite norm fibre and its integral obstructions}

\begin{statement}{A.180}{Proved}
{Stable integer divisors, ordinary ideal classes, and the actual unit obstruction}
Use the fixed fields and norm maps of A.179.1, and write
\(U_n=\mathcal O_{K_n}^{\times}\),
\(U_{n,S}=\mathcal O_{K_n,S_n}^{\times}\), and
\[
 V_n=\ker\bigl(N_{K_n/\Q}:U_n\longrightarrow\{1,-1\}\bigr).
\tag{A.180.1}
\]
Let \(P_{0,n}\) be the distinguished degree-one prime above two,
and put
\[
 \mathfrak R_{q,n}=P_{q,2,n}-2P_{q,0,n},\qquad
 D_n(c)=P_{0,n}+c_{11}\mathfrak R_{11,n}
                         +c_{59}\mathfrak R_{59,n}.
\tag{A.180.2}
\]
The subscripts on the changed primes retain their literal residue
degrees.  Every coherent family of actual field elements with bottom
norm \(-2\) has exactly these divisors for one fixed
\(c=(c_{11},c_{59})\in\Z^2\).

Write \(C_n=\operatorname{Cl}(K_n)\{3\}\).  Its actual generator
\(R_{59,n}=[\mathfrak R_{59,n}]\) gives compatible identifications
and compatible residues
\[
 \begin{gathered}
 C_n=\Z/3^{a_n},\qquad
 [P_{0,n}]_{(3)}=p_nR_{59,n},\qquad
 [\mathfrak R_{11,n}]=r_nR_{59,n},\\
 \mathcal C=\varprojlim_N C_n=\Z_3/I,
 \qquad I=\bigcap_{n\ge1}3^{a_n}\Z_3,
 \qquad p=(p_n),\quad r=(r_n).
 \end{gathered}
\tag{A.180.3}
\]
Here \(C_1=0\), \(a_1=0\), and all the maps in this limit are
the quotient maps induced by ideal norm and the displayed generator.
Let \(\theta_n=[P_{0,n}]_{(3')}\) denote the prime-to-three part
of the full ordinary ideal class.  Define the exact integer profile set
\[
 \mathcal L_{\Z}=
 \{c\in\Z^2:\theta_n=0\text{ for every }n,
                   \ p+c_{11}r+c_{59}=0\text{ in }\mathcal C\}.
\tag{A.180.4}
\]
For every \(c\in\mathcal L_{\Z}\), choose actual generators
\(g_n\) of \(D_n(c)\) with \(N_{K_n/\Q}g_n=-2\).
The element
\[
 \begin{gathered}
 \omega(c)=\left[\left(\frac{g_n}{N_{n+1,n}g_{n+1}}\right)_{n\ge1}\right]
                      \in\varprojlim{}^{1} V_n,\\
 \varprojlim{}^{1}V_n=
 \frac{\prod_{n\ge1}V_n}
 {\{(u_n/N_{n+1,n}u_{n+1})_n:(u_n)_n\in\prod_{n\ge1}V_n\}}
 \end{gathered}
\tag{A.180.5}
\]
is independent of all generator choices.  The natural profile map on
the actual inverse norm fibre has the exact image
\[
 \operatorname{im}\left(
 \{(b_n)\in\varprojlim_N K_n^{\times}:b_0=-2\}
                  \longrightarrow\Z^2\right)
     =\{c\in\mathcal L_{\Z}:\omega(c)=0\}.
\tag{A.180.6}
\]
Each nonempty fibre is a torsor under \(\varprojlim_N V_n\).

The corresponding completed profile map has image precisely
\[
 \mathcal L_{\Z_3}
   =\{c\in\Z_3^2:p+c_{11}r+c_{59}=0\text{ in }\mathcal C\}.
\tag{A.180.7}
\]
This set is nonempty: a lift of \(-p\) in \(\Z_3\), together
with \(c_{11}=0\), supplies an element.  Every completed profile
fibre is a torsor under \(\varprojlim_N\widehat V_{n,3}\).
The exact map comparing its unit obstruction with A.180.5 is
\[
 \varprojlim_N\widehat V_{n,3}\longrightarrow
 \varprojlim_N(\widehat V_{n,3}/V_n)
       \xrightarrow{\partial}\varprojlim{}^{1}_N V_n
       \longrightarrow0.
\tag{A.180.8}
\]
The quotients in this display are abstract abelian groups; the maps
are induced by the actual lattice inclusions and norms.  Thus
A.180.4--6 give an exact test for the discrete actual fibre, including
the full ideal classes and the actual unit compatibility, while
A.180.7 describes all profiles of the completed family of A.179.
No value of \(\theta_n\) for \(n\ge3\), or of \(\omega(c)\),
is asserted by the completed existence theorem.
\end{statement}

\begin{proof}
\emph{The literal divisor maps.}
A.166 proves the equality of the unrestricted and fixed-support
inverse fibres and describes every local divisor coordinate.
At two only the stable degree-one branch survives, with coefficient
one.  At three the coefficient is zero.  At each changed prime the
two residue-degree weights are one and two; the norm-zero vectors
are exactly \((-2c_q,c_q)\).  Every further changed-prime extension
is totally ramified, so the ideal norm weight on either branch is
one.  Thus the coefficients \(c_q\) are constant along the tower.
Actual valuations give integers, and completed valuations give
elements of \(\Z_3\).  In both cases this is exactly the divisor
map A.180.2, with
\[
 N_{n+1,n}D_{n+1}(c)=D_n(c),\qquad
                  N_{n,0}D_n(c)=[2].
\tag{A.180.9}
\]
These formulas also prove uniqueness of the profile.

The actual first-level elements
\[
 \gamma_{11}=2+\alpha,\qquad
 \gamma_{59}=2+2\alpha-\alpha^2,\qquad
 h_q=q/\gamma_q^3
\tag{A.180.10}
\]
have norms \(q,q,1\), respectively, with the appropriate value
of \(q\), as proved by the literal identities in A.165.8.
At level one \(\operatorname{div}(h_q)=\mathfrak R_{q,1}\).
At every changed-prime branch the ramification index from level one
to level \(n\) is \(3^{n-1}\).  Hence in the full ordinary
divisor group one has
\[
                 \operatorname{div}_{K_n}(h_q)
                              =3^{n-1}\mathfrak R_{q,n}.
\tag{A.180.11}
\]
It follows that \([\mathfrak R_{q,n}]\) has three-power order
already in the full ordinary class group.  Its three-primary
projection therefore changes no class.  Consequently the full
class of \(D_n(c)\) is zero exactly when its prime-to-three
component \(\theta_n\) is zero and
\(p_n+c_{11}r_n+c_{59}=0\) in \(\Z/3^{a_n}\).
This proves the asserted full principality test, including the
independence of \(\theta_n\) from the two profile coordinates.

\emph{The compatible class coordinates.}
A.179 identifies \(C_n\) as the cyclic group generated by
\(R_{59,n}\).  The ideal norm of this generator is the preceding
generator, by the residue-degree-one transition on both changed
branches.  Its order exponents are therefore nondecreasing and
the induced map is the literal quotient
\(\Z/3^{a_{n+1}}\to\Z/3^{a_n}\).
Equation A.180.11 gives \(a_n\le n-1\).
The distinguished prime also has norm the preceding distinguished
prime, so the residues \(p_n,r_n\) are compatible.
The map \(\Z_3\to\varprojlim_n\Z/3^{a_n}\) is onto:
the nonempty nested closed sets of lifts of any compatible residue
sequence have a common point in \(\Z_3\).
Its kernel is exactly \(I\).  This proves A.180.3 and identifies
the simultaneous three-primary equations with the single equation
in A.180.4.

For reference the already evaluated levels give
\[
 (a_2,p_2,r_2)=(1,0,2),\qquad
 (a_3,p_3,r_3)=(2,6,5).
\tag{A.180.12}
\]
The latter equation requires
\(c_{59}\equiv4c_{11}+3\pmod9\).
For example \((0,3)\) and \((3,6)\) satisfy every class equation
currently evaluated in A.180.12.  The full simultaneous equation
remains the equation in A.180.4, with its exact restriction maps.

The following useful additional identity concerns the stated
three-primary classes.  If \(P_{f,n}\) is a prime above two of
residue degree \(f\), then
\[
           [P_{f,n}]_{(3)}=f[P_{0,n}]_{(3)},\qquad
                         3^{n-2}p_n=0\text{ in }C_n
                         \quad(n\ge2).
\tag{A.180.13}
\]
At level two the three primes above two are principal, by the
actual elements \(\beta_2,\beta_1/\beta_2,2/(\alpha-1)\)
of A.165.12--13.  At each next level the old distinguished
prime of \(E_{n-1}\) splits into three primes in \(E_n\).
A.177 makes their three-primary classes equal, since the root
automorphism acts trivially on that class group.  Quadratic ideal
extension sends the pure degree-one prime to one of these primes
and the pure degree-two prime to the other two.  This extension
is injective on three-primary groups, because its composite with
norm is multiplication by two.  Thus
\([P_{2,n}]_{(3)}=2[P_{0,n}]_{(3)}\).
All older branches are unramified inert cubic extensions, so
ideal extension takes their preceding prime to their upper prime.
The extension of the preceding distinguished prime is
\(P_{0,n}+P_{2,n}\), whose class is \(3[P_{0,n}]_{(3)}\).
Induction proves the first identity in A.180.13 for every branch.
Finally, the extension of the actual principal divisor
\(\operatorname{div}_{K_2}(\beta_2)=P_{0,2}\) has total residue
degree \(3^{n-2}\) at two.  Applying the first identity to its
unramified prime factors proves the second.

\emph{Actual generators and their precise unit defect.}
For \(c\in\mathcal L_{\Z}\), the proved ordinary class test
makes each \(D_n(c)\) an actual principal fractional ideal.
Choose a generator.  Its rational norm has divisor \([2]\)
by A.180.9 and hence equals either \(2\) or \(-2\).
The degree \([K_n:\Q]=3^n\) is odd, so multiplication by
\(-1\) changes its norm sign.  This produces the prescribed
generator \(g_n\) with norm exactly \(-2\).
The quotient
\(\delta_n=g_n/N_{n+1,n}g_{n+1}\) has zero divisor and
rational norm one.  It is therefore in precisely the group
\(V_n\) of A.180.1.

Replacing \(g_n\) by another such generator \(g_nv_n\)
uses \(v_n\in V_n\) and replaces \(\delta_n\) by
\(\delta_n v_n/Nv_{n+1}\).  This proves the independence
in A.180.5.  If \(\omega(c)=0\), choose actual units \(u_n\)
with \(u_n/Nu_{n+1}=\delta_n\).  Then the literal elements
\[
                              b_n=g_n/u_n
\tag{A.180.14}
\]
have the desired divisors, norm exactly \(-2\), and satisfy
\(N b_{n+1}=b_n\).  Conversely such a family makes
\(u_n=g_n/b_n\) an actual solution of these unit equations.
This proves both directions of A.180.6.  Dividing any two
solutions gives exactly a coherent sequence in \(V_n\),
which proves the torsor assertion.

\emph{The completed divisor sequence and its norm target.}
Let \(\Lambda_n\) be the image of the finite-support divisor
map \(U_{n,S}\to\bigoplus_{P\in S_n^{\rm fin}}\Z[P]\).
The exact sequence
\[
 1\longrightarrow U_n\longrightarrow U_{n,S}
                  \xrightarrow{\operatorname{div}}\Lambda_n
                  \longrightarrow0
\tag{A.180.15}
\]
splits as a sequence of abelian groups because \(\Lambda_n\)
is free.  Its inclusion in the finite free divisor lattice has
finite cokernel, namely the image of that lattice in the ordinary
class group.  Completion therefore gives kernel
\(\widehat U_{n,3}\) and image equal to the kernel of the
divisor lattice's map to the three-primary class group.  One
can obtain this directly by tensoring the finite free lattice
sequence with the flat module \(\Z_3\); the finite cokernel
then retains exactly its three-primary subgroup.

It follows that a completed generator with divisor \(D_n(c)\),
for \(c\in\Z_3^2\), exists exactly when
\(p_n+c_{11}r_n+c_{59}=0\).
Its rational norm is already \([-2]\): the four rational
prime valuations are those of two, and the remaining rational
unit group is \(\{1,-1\}\), whose pro-three completion is
zero.  The quotient of any two such completed generators is
in \(\widehat U_{n,3}\).

The actual field \(K_n\) is contained in \(\R\), so its
roots of unity are \(\{1,-1\}\).  Norm on this subgroup is
the identity, because the degree is odd.  Thus
\[
 U_n=\{1,-1\}\times V_n,
 \qquad V_n\text{ is free of finite rank},\qquad
                \widehat U_{n,3}=\widehat V_{n,3}.
\tag{A.180.16}
\]
The inclusion \(V_n\hookrightarrow\widehat V_{n,3}\)
is injective.  These groups and maps are exactly the unit
groups in A.180.5 and A.180.8.

For any sequence of compact abelian groups with continuous
transition maps, the map
\((u_n)\mapsto(u_n/Nu_{n+1})\) on their product is onto.
Indeed, for any prescribed right sides, each finite collection
of these equations is solved backwards from its last index.
Each equation defines a closed subset of the compact product,
so the finite intersection property solves all the equations.
Apply this to \(\widehat V_{n,3}\) and the defects of any
chosen completed generators.  Their correction gives a
completed coherent family for every profile in A.180.7.
Conversely its divisor gives precisely those class equations.
This proves the completed image and torsor assertions, with
the actual norm maps and rational target retained.

\emph{The exact unit-lattice comparison.}
Put \(Q_n=\widehat V_{n,3}/V_n\), as an abstract quotient.
For a compatible sequence \(q_n\in Q_n\), choose arbitrary
lifts \(\widehat u_n\in\widehat V_{n,3}\).
Compatibility says
\(\widehat u_n/N\widehat u_{n+1}\in V_n\).
Its class in A.180.5 is independent of the lifts and defines
\(\partial(q)\).  The kernel consists exactly of the images
of coherent completed units: a zero defect class can be removed
by actual units, leaving compatible completed lifts.
Finally every sequence of actual units, regarded as a sequence
in the compact completions, has a completed solution of its
defect equations by the preceding compactness argument.
Their cosets form a compatible sequence in \(Q_n\), whose
image under \(\partial\) is the given defect class.
This proves all of A.180.8 without a topological assertion
about the abstract quotients.

In particular, for \(c\in\mathcal L_{\Z}\), choose a
completed solution of
\(\widehat u_n/N\widehat u_{n+1}=\delta_n\).
The compatible cosets \((\widehat u_nV_n)_n\) map to
\(\omega(c)\).  They lift to coherent completed units exactly
when the actual correction A.180.14 exists.
The full class projection in A.180.15 and this connecting map
give the asserted precise comparisons with the completed family.
\end{proof}

The retained certificate
\path{code/two_inert_discrete_fibre_certificate.py}
checks the literal first-field norms, the integral divisor norm
weights, all known level-two and level-three profile congruences,
and the signs in the defect and connecting maps.  The proof above
establishes the infinite assertions.  In particular it does not
replace the actual unit system in A.180.5 by its compact completion.
